\section{The actual automorphic parts and density on the original ideal}
\label{sec:automorphic-installation-v44}

This chapter completes the passage from a nonzero coefficient-valued
automorphic subspace to the full arithmetic part used by the ordinary
Hecke argument.  The representation, its quotient descent, its
polarization, all local conditions, and its spherical polynomial
readings are constructed together.  The last two statements install
these parts in the original \texttt{AutomorphicHeckeSet} and remove
the additional supplier retained in AFS22.

Throughout, the source data are fixed.  Write $F=L^+$ for the maximal
real subfield of the CM field $L$, let $\Omega$ be its specified
algebraic closure, and let $\ell$, $s$, $V$, $Q$, $\mathcal O$,
$\bar r$, and $\iota:\overline{\mathbb Q}_\ell\simeq\mathbb C$
have their original meanings.  Put $n\geq1$, and write
\[
 p=V.\mathrm{pr}:G_L\longrightarrow G_0,\qquad
 G_L=\operatorname{Gal}(\Omega/L),\qquad
 G_F=\operatorname{Gal}(\Omega/F).
\]
The original finite set of places of $F$ is $S'=V.S'$.
The chosen integral type is $\tau_0$ and satisfies the literal
\texttt{TypeForAlt} condition.  All original Teichm\"uller families,
all correspondence maps satisfying \texttt{IsRecT}, and all coefficient
and period bridges remain quantified as in the source.  The fixed
embedding $\iota$ is used algebraically; its continuity is never assumed.

The inherited hypotheses are the source's
\texttt{HypothesesH'} and its stated literature and period foundations.
The local-transport claims LTP5--LTP21 and Galois-assembly claim GUA18 supply their indicated
consumers.  No additional \texttt{PartAltAlt},
\texttt{AutomorphicHeckeSet}, \texttt{AtLT}, or density assumption is
introduced here.  In particular, the original base-change and
cuspidal Galois-attachment foundations remain identified at the scope
recorded in those chapters.

\subsection{The original quotient and its pairing}

% BEGIN UPA1
\paragraph{UPA1. The place and ramification consequences of the Hida setting.}
\label{claim:UPA1}
The prescribed extension $L/F$ is unramified at every finite place.
For every chosen place $s.w(v)$ one has
$\underline{s.w(v)}=v$, where an underline means restriction to $F$.
Moreover
\[
 \ell\in w\quad\Longrightarrow\quad \underline w\in S',
 \qquad
 \underline w\in S'\quad\Longrightarrow\quad w\text{ is split over }F.
 \tag{UPA.1}
\]
These conclusions apply simultaneously at the original $\ell$-adic,
type, auxiliary, and distinguished places.
\begin{proof}
The condition $s.\mathrm{unr}$ says that every absolute inertia subgroup
over a finite place of $F$ is contained in the image of
$G_L\hookrightarrow G_F$.  The image of absolute inertia in
$\operatorname{Gal}(L/F)$ is the finite inertia subgroup.  It is
therefore trivial.  Residue fields of number fields are perfect, so
the usual equality of the inertia order with the ramification index
gives ramification index one at every finite place.  This is exactly
the finite-place unramifiedness used by the Galois-assembly theorem.

The valuation ring $s.A(v)$ lies over both $v$ and $s.w(v)$.
Pulling its maximal ideal back through
$\mathcal O_F\longrightarrow\mathcal O_L$ identifies the restriction
of $s.w(v)$ with $v$.  The same pullback identity gives
$\ell\in w$ if and only if $\ell\in\underline w$.
All places of $F$ over $\ell$ belong to $s.S$, and $s.S\subset S'$.
This proves the first implication in (UPA.1).

For $v\in S'$, the original level condition puts the chosen
decomposition group at $v$ in the image of $G_L$.
The image of that group in the quadratic Galois group
$\operatorname{Gal}(L/F)$ is consequently trivial.  All prolongations
over $v$ have conjugate decomposition groups, and the image of $G_L$
is normal.  The decomposition group in the quadratic extension is
therefore trivial at every place above $v$.  Its orbit on the places
above $v$ has two elements; this is the literal \texttt{IsSplitAt}
condition.  The definitions of $s.S$ and $Q$ place the original type
places, the distinguished place, and the auxiliary places in $S'$.
The auxiliary places are away from $\ell$ by the already checked
original \texttt{axU\_s} conditions.
\end{proof}
% END UPA1

% BEGIN UPA2
\paragraph{UPA2. Inertia outside the level kills the actual quotient kernel.}
\label{claim:UPA2}
Let $K\leq G_L$ be a closed subgroup.  Suppose that it contains every
inertia subgroup at every prolongation of every place $w$ of $L$
with $\underline w\notin S'$.  Then
\[
                         \ker p\subset K.                    \tag{UPA.2}
\]
\begin{proof}
First record the valuation-theoretic point used in passing between
the two number fields.  If a valuation ring $A\subset\Omega$ lies
over a finite place $v$ of $F$, then every element of $\mathcal O_L$
is integral over $A$ and hence lies in $A$.  The inverse image of
the maximal ideal of $A$ in $\mathcal O_L$ is a prime $w$ lying
over $v$.  It is nonzero because its contraction is $v\ne0$.
Thus it is a finite place of $L$, and $A$ lies over $w$.

Write $H$ for the image of $G_L$ in $G_F$, and $N_{S'}$ for the
closed normal subgroup generated by the absolute inertia groups
at the places of $F$ outside $S'$.  This is the original subgroup
used to define the level quotient.  The image of $K$ in $G_F$
is closed: $K$ is compact, restriction is continuous, and $G_F$
is Hausdorff.

Consider a generator $g$ from inertia at $A$ over
$v\notin S'$ and any conjugate $bgb^{-1}$ in $G_F$.
It belongs to inertia at the transported valuation ring
$A\circ b^{-1}$ over the same place $v$.  The condition
$s.\mathrm{unr}$ places this inertia group in $H$.
Choose the place $w$ of $L$ under $A\circ b^{-1}$.
For an automorphism fixing $L$, inertia over $F$ at that
valuation ring is the same inertia condition as inertia over $L$.
The lift of $bgb^{-1}$ to $G_L$ therefore belongs to $K$,
because $\underline w=v\notin S'$.
Every conjugate of every generator of $N_{S'}$ belongs to
the image of $K$.  Closedness now gives
$N_{S'}\subset\operatorname{res}(K)$.

If $t\in\ker p$, its restriction to $F$ belongs to $N_{S'}$,
hence equals the restriction of some $k\in K$.
Restriction $G_L\hookrightarrow G_F$ is injective on the
underlying automorphisms of $\Omega$, so $t=k\in K$.
This proves (UPA.2) for the original quotient, without replacing
its defining normal closure.
\end{proof}
% END UPA2

% BEGIN UPA3
\paragraph{UPA3. Continuous descent preserves the same polarization matrix.}
\label{claim:UPA3}
Let $R:G_L\to\operatorname{GL}_n(\overline{\mathbb Q}_\ell)$
be continuous and trivial on inertia outside $S'$.
There is a unique continuous representation
$\rho:G_0\to\operatorname{GL}_n(\overline{\mathbb Q}_\ell)$
with $\rho\circ p=R$.
For every matrix $P$ giving the original pairing
\[
 R(g)^{\mathsf t}P R(cgc^{-1})
                =\varepsilon_\ell(g)^{\,1-n}P,
 \tag{UPA.3}
\]
the same $P$ gives the original \texttt{IsForm} for $\rho$.
\begin{proof}
The map $p$ is surjective.  Indeed, lift an element of $G_0$
to an element of $G_F$ in the quotient by $N_{S'}$.
The level identity $V.h\eta$ identifies the inverse image of
$G_0$ with the image of $G_L$, so that lift comes from $G_L$.
This proves surjectivity for the actual map $p$.

The kernel of $R$ is closed.  It is the finite intersection
of the closed conditions $R(g)_{ij}=\delta_{ij}$, since
$\overline{\mathbb Q}_\ell$ is Hausdorff and each matrix entry
is continuous.  UPA2 gives $\ker p\subset\ker R$.
Define $\rho(p(g))=R(g)$.  The kernel inclusion proves
well-definedness, and surjectivity proves uniqueness and the
homomorphism law.

The group $G_L$ is compact.  The target $G_0$ is Hausdorff
because $N_{S'}$ is closed, and $p$ is continuous by the
original restriction and quotient construction.  Thus $p$
is a closed surjection and a quotient map.  The identity
$\rho\circ p=R$ proves continuity of $\rho$.

The frame's involution and cyclotomic character satisfy,
by their original defining identities,
\[
 \operatorname{conj}(p(g))=p(cgc^{-1}),\qquad
 \varepsilon_{G_0}(p(g))=\varepsilon_\ell(g).
 \tag{UPA.4}
\]
Here $cgc^{-1}$ denotes the unique automorphism fixing $L$
whose restriction to $F$ is the displayed conjugate; it
exists because $L/F$ is normal.  In coordinates over
$\mathcal O$, the second identity is precisely the
$V.h\varepsilon$ identity followed by the original scalar
map to $\overline{\mathbb Q}_\ell$.
Substitution of (UPA.4) and $\rho(p(g))=R(g)$ into
(UPA.3), followed by surjectivity of $p$, gives
\[
 \rho(h)^{\mathsf t}P\rho(\operatorname{conj}h)
                   =\varepsilon_{G_0}(h)^{\,1-n}P
 \quad(h\in G_0).
 \tag{UPA.5}
\]
Invertibility of $P$ is unchanged.  This is the literal
source \texttt{IsForm}, with no new pairing or scalar extension.
\end{proof}
% END UPA3

% BEGIN UPA4
\paragraph{UPA4. Semisimplicity and irreducibility pass through the quotient.}
\label{claim:UPA4}
For a surjection $p:G_L\to G_0$ and $R=\rho\circ p$,
$R$ and $\rho$ have exactly the same invariant subspaces.
Consequently $R$ is semisimple if and only if $\rho$ is
semisimple, and $R$ is irreducible if and only if $\rho$
is irreducible.
\begin{proof}
A $\rho$-invariant subspace is $R$-invariant by composition.
Conversely, for $h\in G_0$ choose $g$ with $p(g)=h$;
$R$-invariance then gives invariance under $\rho(h)$.
The same argument applies to any proposed invariant
complement.  The definitions of semisimplicity and
irreducibility now give the two equivalences.  The
underlying vector space and all of its subspaces are
unchanged.
\end{proof}
% END UPA4

\subsection{Spherical polynomials with the original Frobenius convention}

% BEGIN UPA5
\paragraph{UPA5. The spherical Satake theorem has the required normalization.}
\label{claim:UPA5}
For $k=L_w$, $K=\operatorname{GL}_n(\mathcal O_k)$,
and a smooth irreducible complex representation $\pi$ with
$\pi^K\ne0$, there is a multiset
$B=\{\beta_1,\ldots,\beta_n\}\subset\mathbb C^\times$
such that
\[
 {\bf1}_{K\operatorname{diag}(\varpi I_j,I_{n-j})K}
       \ \text{acts on }\pi^K\text{ by }\
 q^{j(n-j)/2}e_j(B)\quad(0\leq j\leq n).
 \tag{UPA.6}
\]
Here $\operatorname{vol}(K)=1$ and $q=\#(\mathcal O_k/\mathfrak m_k)$.
\begin{proof}[Published theorem and exact application]
Apply the spherical Satake isomorphism to the split reductive group
$\operatorname{GL}_n/k$.  A precise reference is B.~H.~Gross,
\emph{On the Satake isomorphism}, Proposition~3.6,
formula~(3.14), and Proposition~6.4.\footnote{
\url{https://people.math.harvard.edu/~gross/preprints/sat.pdf},
pp.~6, 8, 11; the Haar normalization is (2.5).}
The hyperspecial group is exactly $K$ and the Haar normalization
is the one in (UPA.6).  The cocharacter $(1^j,0^{n-j})$
is minuscule; its dual representation is $\bigwedge^j\mathbb C^n$,
whose character is $e_j$.  Its pairing with the half-sum of
positive roots is $j(n-j)/2$.  Thus the stated formula is
the normalized Satake formula.  The parameter lies in
$\operatorname{GL}_n(\mathbb C)$, so all $\beta_i$ are nonzero.
The result holds for every finite extension of $\mathbb Q_p$
and every $n\geq1$; the source places meet these hypotheses.
The fixed space is nonzero and one dimensional, the multiset has
cardinality $n$, and (UPA.6) holds on every fixed vector.
Changing the uniformizer by a unit leaves the double cosets unchanged.
Thus $B$ satisfies the literal source $\mathrm{HasSatGL}$ predicate,
with all its uniformizer and fixed-vector quantifiers; it is exactly
the normalized parameter used in GUA12.
\end{proof}
% END UPA5

% BEGIN UPA6
\paragraph{UPA6. The original spherical operators give the original Hecke polynomial.}
\label{claim:UPA6}
For the representation and multiset in UPA5, put
\[
 a_j=q^{h(n-h)/2}e_h(B),\qquad h=\min(j,n).
 \tag{UPA.7}
\]
Then $a_0=1$, every original operator $\operatorname{sphOp}_j$
acts on $\pi^K$ by $a_j$, and
\[
 \prod_{i=1}^n\bigl(X-q^{(n-1)/2}\beta_i\bigr)
 =\sum_{j=0}^n(-1)^j q^{j(j-1)/2}a_jX^{n-j}
 =\operatorname{heckePol}(n,q,a).
 \tag{UPA.8}
\]
\begin{proof}
The source's compact subgroup $K_O(k,n)$ equals $K$:
both consist of invertible integral matrices with integral inverse.
Its chosen uniformizer has valuation one for the normalized
discrete valuation.  Any two such choices differ by a unit,
and therefore determine the same diagonal double coset.

The source matrix $\operatorname{al}(k,n,j)$ has a uniformizer
in its first $\min(j,n)$ diagonal entries and one elsewhere.
For a $K$-fixed vector, its double-coset operator is the
sum of $\pi(x)$ over the right cosets $xK$ in that double
coset.  This is exactly convolution by its characteristic
function when $\operatorname{vol}(K)=1$; compactness makes
the sum finite.  UPA5 proves (UPA.7) for all $j$, including
$j>n$.  At $j=0$ the operator and its eigenvalue are one.

Expand the product on the left of (UPA.8) by elementary
symmetric functions.  Its coefficient of $X^{n-j}$ is
$(-1)^j q^{j(n-1)/2}e_j(B)$.  Since
\[
       \frac{j(n-1)}2
          =\frac{j(j-1)}2+\frac{j(n-j)}2,
\]
this is the coefficient on the right.  The number
$j(j-1)/2$ is an integer, and all real powers here use the
positive real $q$.  Thus (UPA.8) is an equality with the
literal source polynomial and its original normalizations.
\end{proof}
% END UPA6

% BEGIN UPA7
\paragraph{UPA7. Spherical local compatibility gives the geometric Frobenius polynomial.}
\label{claim:UPA7}
Let $w\nmid\ell$, let $\pi$ be spherical with Satake multiset
$B$, and let $R$ be a continuous global representation.
Suppose that, for every prolongation $A$ of $w$ and every
$\mathrm{rc}$ satisfying \texttt{IsRecT},
\[
 \texttt{WDT}\bigl(\iota,R,\mathrm{rc}(\operatorname{algNorm}\pi)\bigr)
 \tag{UPA.9}
\]
holds.  Then $R$ is unramified at $w$ and, for every
arithmetic Frobenius lift $g$ at $A$,
\[
 \iota\!\left(\operatorname{charpoly}R(g^{-1})\right)
       =\prod_{\beta\in B}(X-q^{(n-1)/2}\beta)
       =\operatorname{heckePol}(n,q,a).
 \tag{UPA.10}
\]
\begin{proof}
Choose a correspondence using the retained nonemptiness theorem
\texttt{RecTExists}.  The spherical correspondence calculation
GUA12, with the geometric reciprocity
normalization, gives
\[
 \mathrm{rc}(\operatorname{algNorm}\pi)
   =\left(\bigoplus_{\beta\in B}\operatorname{unr}
       \bigl(\beta^{-1}q^{(1-n)/2}\bigr),\,0\right).
 \tag{UPA.11}
\]
This calculation follows also by ordering the unramified
rank-one characters, taking their Langlands quotient, and
applying the \texttt{one}, \texttt{langlands}, and
\texttt{nu} clauses of \texttt{IsRecT}; no global uniqueness
assertion for the weaker correspondence is used.

The parameter in (UPA.11) has zero monodromy and is trivial
on commutators of the Weil group, because every summand is
a character.  It therefore satisfies the source predicate
\texttt{Gd}.  Unpacking \texttt{WDT}, its witnessing
Weil--Deligne parameter must equal (UPA.11).
The zero-monodromy argument of UPA9, applied with trivial inertia
characters, makes the original inertia action trivial: an open
inertia subgroup is killed, every inertia matrix consequently has
finite order, and the commuting semisimple-plus-nilpotent comparison
then identifies it with the identity.  For every Weil element, the
same commuting semisimple-plus-nilpotent decomposition preserves
its characteristic polynomial.  This proves both unramifiedness
and the required polynomial comparison.  Arithmetic Frobenius
has degree one.  Its inverse has degree minus one, so its
eigenvalues in (UPA.11) are
\[
       \bigl(\beta^{-1}q^{(1-n)/2}\bigr)^{-1}
                          =q^{(n-1)/2}\beta.
\]
This proves the first equality of (UPA.10); UPA6 proves
the second.  In particular the inverse on $g$ is part of
the conclusion and has not been absorbed by renaming the
Satake parameters.
\end{proof}
% END UPA7

\subsection{Type and ordinary conditions for the same representation}

% BEGIN UPA8
\paragraph{UPA8. A global-character inertia type activates the exact WD comparison.}
\label{claim:UPA8}
Suppose a Weil--Deligne class $y$ has zero monodromy and
diagonal inertia characters that are restrictions of
characters $\chi_i:G_L\to\mathbb C^\times$.
Then $y$ satisfies \texttt{Gd}.  Consequently
\texttt{WDT}$(\iota,R,y)$ implies
\texttt{IsWDOf}$(\iota,R,y)$.
\begin{proof}
Choose the representative occurring in the literal
\texttt{IsTypeWD} condition.  Every Weil commutator has
degree zero and hence belongs to inertia.  Each $\chi_i$
is trivial on it, because its target is commutative.
The representative therefore acts as the identity on
every commutator.  Its characteristic polynomial there
is $(X-1)^n$, which is precisely the source torality
condition \texttt{Tor}.  Its monodromy is zero by
hypothesis.  These two assertions are \texttt{Gd}.
The definition of \texttt{WDT} provides a parameter $x$
with \texttt{IsWDOf}$(\iota,R,x)$ and an implication
$\texttt{Gd}(y)\Rightarrow x=y$; substitution gives the
last assertion.
\end{proof}
% END UPA8

% BEGIN UPA9
\paragraph{UPA9. Zero-monodromy type gives literal diagonal inertia over the original field.}
\label{claim:UPA9}
Let $w\nmid\ell$, let $A$ lie over $w$, and suppose
\texttt{IsWDOf}$(\iota,R,y)$.
If $y$ has zero monodromy and inertia type
$\bigoplus_{i=0}^{n-1}\iota\chi_i$, with
$\chi_i:I_A\to\overline{\mathbb Q}_\ell^\times$,
then there is $C\in\operatorname{GL}_n(\overline{\mathbb Q}_\ell)$
such that, simultaneously for all $t\in I_A$,
\[
               C R(t) C^{-1}=\operatorname{diag}(\chi_i(t)).
 \tag{UPA.12}
\]
\begin{proof}
Unpack the source \texttt{IsWDOf} witness.  It contains a
monodromy matrix $N$ over $\overline{\mathbb Q}_\ell$, an
open inertia subgroup $U$ on which $R$ is the exponential
of a scalar multiple of $N$, and a Frobenius-semisimple
complex representative $(r_{\mathrm{WD}},N_{\mathrm{WD}})$.
The given type identifies this representative with one
having zero monodromy.  The monodromy intertwining
identity implies $N_{\mathrm{WD}}=0$; injectivity of the
matrix scalar map induced by $\iota$ gives $N=0$.
Thus $R|_U=1$.

For every $t\in I_A$, some positive power belongs to $U$.
Indeed, $U$ is open in the compact inertia group, hence
has finite index; two powers of $t$ have the same coset.
It follows that $R(t)$ has finite order.  Over a
characteristic-zero field its minimal polynomial divides
$X^m-1$ for some $m>0$, a separable polynomial.
Thus $\iota R(t)$ is semisimple.

The source Frobenius-semisimplification witness says that
$\iota R(t)$ commutes with $r_{\mathrm{WD}}(t)$ and that
their difference is nilpotent.  Here the usual monodromy
exponential is one because $N=0$.
The second operator is also semisimple.  Commuting
semisimple operators over $\mathbb C$ are simultaneously
diagonalizable, so their difference is semisimple.
A semisimple nilpotent operator is zero.  Hence
\[
                   \iota R(t)=r_{\mathrm{WD}}(t)
                   \quad(t\in I_A).                         \tag{UPA.13}
\]
The fixed intertwiner supplied by the type conjugates
the right-hand side of (UPA.13) to
$\operatorname{diag}(\iota\chi_i(t))$ for all $t$ at once.
Apply $\iota^{-1}$ entrywise to that intertwiner and its
inverse.  Their product identities remain valid and give
the matrix $C$ in (UPA.12).
This is algebraic transport of matrices; no topological
property of $\iota$ is required.  Repeated characters cause
no difficulty.
\end{proof}
% END UPA9

% BEGIN UPA10
\paragraph{UPA10. Distinct-unit principal series satisfy the exact plain-toral predicate.}
\label{claim:UPA10}
Let $\pi\simeq\operatorname{PS}(\omega)$ be smooth and
irreducible, where all $\omega_i:k^\times\to\mathbb C^\times$
are smooth and their restrictions to $\mathcal O_k^\times$
are pairwise distinct.  Then the irreducible class of
$\pi$ satisfies the source predicate \texttt{PlainTor}.
\begin{proof}
By the definition of $\operatorname{PS}$, its normalized
induction characters are
\[
              \mu_i=\omega_i|\cdot|^{(n-1)/2}.
 \tag{UPA.14}
\]
The original cuspidal-support theorem identifies the
cuspidal support of the irreducible principal series with
the multiset of these rank-one characters.  Thus every
support member has rank one, the first clause of
\texttt{PlainTor}.

No two support members differ by the source unit norm
twist.  For distinct indices, restrict a proposed equality
$\mu_i=\mu_j|\cdot|$ to units.  Both norm factors become
one, contradicting the distinct-unit hypothesis.
For the same index it would give $|\cdot|=1$, contradicted
at a uniformizer because $q>1$.  This is the second
clause of \texttt{PlainTor}.  The support assertion uses
the original local classification foundation already
retained in the principal-series arguments; no assertion
about arbitrary nontoral principal series is needed.
\end{proof}
% END UPA10

% BEGIN UPA11
\paragraph{UPA11. The infinity multiset yields the original Hodge--Tate integers.}
\label{claim:UPA11}
Suppose the original multiset identity at a local embedding
$j:L_w\to\overline{\mathbb Q}_\ell$ is
\[
 \iota_{\mathbb Z}\bigl(\mathrm{HT}(j|_L)\bigr)
   =\left\{\,\frac{n-1}{2}-z:
                    z\in\operatorname{infOfWeight}(\lambda(j|_L)^\vee)
     \,\right\}.
 \tag{UPA.15}
\]
Then
\[
 \mathrm{HT}(j|_L)
       =\{\,\lambda(j|_L)_{n-1-i}+i:0\leq i<n\,\}.
 \tag{UPA.16}
\]
Every embedding arising from an original local period bridge
is a restriction $j|_L$ of this kind.
\begin{proof}
The source conventions give
$\lambda^\vee_i=-\lambda_{n-1-i}$ and
\[
 \operatorname{infOfWeight}(\lambda^\vee)
   =\left\{-\lambda_{n-1-i}+\frac{n-1}{2}-i:0\leq i<n\right\}.
 \tag{UPA.17}
\]
Subtracting each element from $(n-1)/2$ gives the integers
in (UPA.16), with their multiplicities.  The map
$\mathbb Z\hookrightarrow\mathbb C$ is injective, so its
map on finite multisets is injective.  This proves
(UPA.16).  The displayed integers are exactly the
source functions $\operatorname{htW}(\lambda,i,j|_L)$.

For the assertion about a period bridge, let $\tau:L\to
\overline{\mathbb Q}_\ell$ be induced by the bridge and
an embedding of its finite local coefficient field.
Embeddings between finite extensions of $\mathbb Q_\ell$
preserve the extended valuation: the absolute value is
the unique extension of the one on $\mathbb Q_\ell$,
equivalently the spectral norm determined by the minimal
polynomial.  The bridge's valuation identity therefore
shows that the finite place induced by $\tau$ is exactly
$w$.  The map $\tau$ extends continuously to the completion
$L_w$.  Its image lies in the finite complete extension
generated by $\tau(L)$ and $\mathbb Q_\ell$, so this
extension still takes values in $\overline{\mathbb Q}_\ell$.
It is the required $j$, and $j|_L=\tau$.
\end{proof}
% END UPA11

% BEGIN UPA12
\paragraph{UPA12. The original local-global and Hodge data imply the full AtLT condition.}
\label{claim:UPA12}
Suppose $R$, $\mathrm{HT}$, $w$, and $\pi$ satisfy
\texttt{LocGlobT}; suppose
\texttt{PlainTor}$(\pi)$ implies
\texttt{HasHTAt}$(R,\mathrm{HT},w)$; and suppose the
infinity-weight identity (UPA.15) holds at every local
embedding.  Then the original
\texttt{AtLT}$(\iota,A,R|_{D_A},\lambda,\pi)$ holds.
\begin{proof}
Unpack an arbitrary test in the definition of
\texttt{AtLT}.  It gives smooth characters $\omega_i$
with pairwise distinct unit restrictions and an
equivalence $\pi\simeq\operatorname{PS}(\omega)$.
By UPA10, $\pi$ is plain toral.  The stated Hodge input
therefore gives, for every original period bridge,
the de Rham assertion and its precise jump multiset.
UPA11 identifies that multiset with
$\{\operatorname{htW}(\lambda,i,\tau)\}_i$ for every
embedding induced by that bridge.

Fix also the test's arbitrary $\mathrm{rc}$ satisfying
\texttt{IsRecT}.  The plain-toral parameter calculation in
GUA2 gives
\[
           \texttt{Gd}\bigl(\mathrm{rc}
                      (\operatorname{algNorm}\pi)\bigr).
 \tag{UPA.18}
\]
Applying the second clause of \texttt{LocGlobT} to each
of the same period bridges now gives potential
semistability and equality of its Weil--Deligne class
with $\mathrm{rc}(\operatorname{algNorm}\pi)$.
In particular their characteristic polynomials agree
on every Weil element.

The de Rham assertion, the jump assertion, and this
characteristic-polynomial identity are exactly the
three clauses of the source \texttt{DRat} predicate.
They hold for every test and every bridge, so they
give \texttt{AtLT}.  The additional \texttt{AbW}
hypothesis in each test is retained and causes no
restriction in this argument.  No ordinary flag has
been silently substituted for the stated predicate.
\end{proof}
% END UPA12

\subsection{The full part and its point properties}

% BEGIN UPA13
\paragraph{UPA13. One representation and one family of local components suffice.}
\label{claim:UPA13}
Fix an original algebraic weight $\lambda$, its actual basis
$e$ and representation $\xi$, and an automorphic irreducible
subspace $\sigma$ for which
\[
                         Y_\sigma=YOf(Q,\iota,\tau_0,e,\xi,\sigma)
                                      \ne0.
 \tag{UPA.19}
\]
There are a continuous semisimple representation
$\rho:G_0\to\operatorname{GL}_n(\overline{\mathbb Q}_\ell)$,
one multiset function $\mathrm{HT}$, and a choice
$\pi_w$ of local component at every split place, such that:
the pullback $R=\rho\circ p$ has all the conclusions of
the literal \texttt{GalUnitaryT} theorem; every $\pi_w$
is a component of the same $\sigma$; and $\pi_w$ is spherical
whenever $\sigma$ is unramified at $w$.
\begin{proof}
Use the split-component existence and fixedness statement
from LTP19--LTP21, choosing $\pi_w$ for
each split $w$.  Its fixed-vector assertion applies to
these very components, so later unramifiedness does not
require a second component choice.

UPA1 verifies the finite-place unramifiedness hypothesis
for the CM extension.  The setting's original odd-prime
condition verifies the other prime hypothesis.  Apply
GUA18 to the original $\sigma$,
rank $n$, and fixed $\iota$.  It produces a single $R$
and $\mathrm{HT}$ with continuity, semisimplicity,
the stated inertia vanishing, the weight identity,
local compatibility at every split place and for every
component, the pairing, and the conditional cuspidal
conclusion.

LTP12 proves that $\sigma$ is
unramified at every $w$ with $\underline w\notin S'$.
By UPA1 these places are away from $\ell$.  Thus the
inertia-vanishing clause of GUA14 and GUA18
applies at every prolongation of each such $w$.
UPA3 descends $R$ to $\rho$ through the actual $p$,
and UPA4 gives semisimplicity of $\rho$.
All the stated local and weight identities are
identities for the same pullback $R=\rho\circ p$.
\end{proof}
% END UPA13

% BEGIN UPA14
\paragraph{UPA14. The sixteen fields of the original arithmetic part are filled.}
\label{claim:UPA14}
Under the data of UPA13 there exists an original
\texttt{PartAltAlt} $Y$ with $Y.Y=Y_\sigma$.
Its representation is the $\rho$ of UPA13, and its local
components are the restrictions of the one chosen family
$\pi_w$ to the original $\ell$-adic and auxiliary places.
\begin{proof}
We specify every field.
Set $Y.Y=Y_\sigma$.  The definition of $YOf$ is an
intersection with the original \texttt{Big} space, so
$Y.Y\leq\texttt{Big}$.  Nonzeroness is (UPA.19).
At an $\ell$-adic index $v$, set
$Y.\sigma_l(v)=\pi_{s.w(v)}$; at an auxiliary index $q$,
set $Y.\sigma_q(q)=\pi_{Q.wq(q)}$.
UPA1 proves that these places are split, and
LTP7 and LTP9 give the exact
\texttt{isoL} and \texttt{isoQ} fields.

For each original spherical index $w\in H.Wp$, its
underlying place is split and lies outside $S'$.
The component $\pi_w$ chosen in UPA13 has nonzero
maximal-compact fixed vectors.  Apply UPA5 and UPA6
to obtain $B_w$ and the full sequence $a(w,j)$.
This gives $a(w,0)=1$.  LTP11
transfers the scalar action of every
$\operatorname{sphOp}_j$ from $\pi_w$ to every element
of $Y_\sigma$.  These are the \texttt{a},
\texttt{a\_zero}, and \texttt{sph} fields.

Set $Y.r=\rho$.  Its continuity and semisimplicity
are UPA13.  At each spherical index, the away-from-$\ell$
clause of \texttt{LocGlobT} gives (UPA.9) for every
prolongation and every permitted correspondence.
UPA7 therefore gives, for every original Frobenius lift $g$,
\[
 \iota\!\left(\operatorname{charpoly}
                  \rho\bigl(p(g^{-1})\bigr)\right)
       =\operatorname{heckePol}(n,\operatorname{Norm}w,a(w,-)).
 \tag{UPA.20}
\]
This is \texttt{frob}, including its specified inverse.

At each $\ell$-adic place LTP14
identifies the infinity weight of $\sigma$
with the original dual of $\lambda$.  The weight clause
of GUA15 and GUA18 is therefore (UPA.15).
Its local clause gives \texttt{LocGlobT} and the
plain-toral implication to \texttt{HasHTAt}.
UPA12 supplies the entire \texttt{atL} field for
the original local representation
$\rho\circ H.\mathrm{dec}(v)=R|_{D_{s.A(v)}}$.

Each auxiliary place is away from $\ell$ by the original
$Q$-conditions.  Taking the first clause of
\texttt{LocGlobT} gives \texttt{WOT} for every
prolongation and every \texttt{IsRecT} map, exactly
the \texttt{atQ} field.
We have supplied
\[
\begin{gathered}
Y,\ \mathrm{le},\ \mathrm{ne},\
\sigma_l,\ \mathrm{isoL},\
\sigma_q,\ \mathrm{isoQ},\
a,\ a_{\mathrm{zero}},\ \mathrm{sph},\\
r,\ \mathrm{cont},\ \mathrm{ss},\
\mathrm{frob},\ \mathrm{atL},\ \mathrm{atQ}.
\end{gathered}
\]
Every entry uses the original function space and the same $\rho$.
\end{proof}
% END UPA14

% BEGIN UPA15
\paragraph{UPA15. The same part satisfies all four unitary point properties.}
\label{claim:UPA15}
The part $Y$ of UPA14 satisfies the original
\texttt{UnitaryPointProperties}$(\bar r,\iota,\lambda,Y.r)$:
its \texttt{form}, \texttt{type}, \texttt{cusp}, and
\texttt{aut} fields all hold with the original quantifiers.
\begin{proof}
For \texttt{form}, apply the pairing clause of
GUA16 to the setting's chosen complex
conjugation.  UPA3 transports its same invertible matrix
to the original frame, with multiplier
$\varepsilon_{G_0}^{1-n}$.

For \texttt{type}, fix an arbitrary Teichm\"uller family
$t$ for $\bar r$ and an arbitrary original type index $u$.
The condition \texttt{TypeForAlt}, together with the
nonzero tensor-intertwiner constructed in
LTP17--LTP18, proves the literal
\texttt{IsTypeWD} condition for $\pi_{s.w(u)}$ and
every allowed correspondence.  Choose one such
correspondence using \texttt{RecTExists}.  The characters
on inertia are the restrictions of the actual global
characters
\[
 G_L\xrightarrow{p}G_0\xrightarrow{t_i}\mathcal O^\times
       \longrightarrow\overline{\mathbb Q}_\ell^\times.
 \tag{UPA.21}
\]
The type place is away from $\ell$.  Local compatibility
there gives \texttt{WDT}, and UPA8 upgrades it to
\texttt{IsWDOf}.  UPA9 supplies one matrix $C$ conjugating
the full inertia action to the ordered diagonal of (UPA.21).
The frame subgroup $I_u$ is the image of that actual
inertia group under $p$.  Taking preimages proves
\[
 C\rho(x)C^{-1}
       =\operatorname{diag}\bigl(\mathcal O\to
                       \overline{\mathbb Q}_\ell(t_i(x))\bigr)
       \quad(x\in I_u).
 \tag{UPA.22}
\]
The choices of $t$ and $u$ were arbitrary.  In particular
the assertion retains every Teichm\"uller family, not
only a family used earlier to construct $\tau_0$.

For \texttt{cusp}, suppose $\rho$ is irreducible.
UPA4 makes $R=\rho\circ p$ irreducible.
GUA17 gives a regular algebraic,
conjugate-self-dual cuspidal $\pi$ attached to $R$.
Its infinity-character clause, applied to the dual-weight
identities of LTP14, gives the original
$\operatorname{HasWt}(\iota,s,\pi,\lambda)$.
This proves precisely the conditional cusp assertion.

For \texttt{aut}, use the original $\sigma$ as witness.
It is automorphic by hypothesis.  Its required
all-component \texttt{LocGlobT} is supplied by
GUA15 and GUA18 for the same $R$.  LTP12 and LTP14
give its dual infinity weights and unramifiedness
outside $S'$.  The family in UPA13 gives a local component
at every split place.  At $s.w(s.v_0)$ use that very
component; LTP15
gives its nonzero $K_c(1)$-fixed vector.
These are every clause in the source \texttt{aut} field.
\end{proof}
% END UPA15

% BEGIN UPA16
\paragraph{UPA16. The strengthened PartX statement and the original PartU follow.}
\label{claim:UPA16}
The constructions above prove the full original
\texttt{partX\_alt}, and hence \texttt{partU\_alt}, on
their original inputs.
In addition to UPA14--UPA15, the chosen $\ell$-adic components
are actual components of $\sigma$ and
\texttt{LocGlobT} holds for every split place and every
component of $\sigma$.  If $Y.r\circ p$ is irreducible,
the resulting cusp is regular algebraic and conjugate
self-dual, is attached to $Y.r\circ p$, is spherical
wherever $\sigma$ is unramified, and has the required
split local-component agreement.
\begin{proof}
The two local-component assertions were retained
explicitly in UPA13 and were used unchanged in UPA14.
For the strengthened final implication, apply the
full conditional cuspidal clause of the same
GUA17--GUA18 to $R=Y.r\circ p$.
It gives all the listed assertions for one cusp,
including its Satake multiset at every place where
$\sigma$ is unramified and the implication from each
of its split local components to a component of
$\sigma$.  UPA15 supplies the four point properties,
and UPA14 supplies $Y.Y=Y_\sigma$.  These form the
entire tuple required by \texttt{partX\_alt}.
Projecting away its strengthened final implication
gives exactly \texttt{partU\_alt}.
\end{proof}
% END UPA16

\subsection{Installation and the original same-ideal density theorem}

% BEGIN UPA17
\paragraph{UPA17. The original AutomorphicHeckeSet is constructed.}
\label{claim:UPA17}
For every original $\tau_0$ satisfying \texttt{TypeForAlt},
\[
 (hset(Q,\tau_0)).\operatorname{AutomorphicHeckeSet}
       \bigl(\iota,\operatorname{UnitaryPointProperties}
                              (\bar r,\iota)\bigr)
 \tag{UPA.23}
\]
holds under the original \texttt{HypothesesH'}.
\begin{proof}
Fix an arbitrary admitted weight, an arbitrary algebraic
coefficient representation $\xi$, and a function $F$
meeting the source's \texttt{Big}, away-level, and
$I_1(e)$ fixedness conditions.  Choose the actual
algebraic-weight basis $e$ provided by
\texttt{IsAlgWt}; the choices of level and depth
are kept as supplied.

ITA2 gives a Teichm\"uller family.  Apply
\texttt{TypeForAlt} to that family to see that the
complexified type has open kernel.  The scalar map
$\mathcal O\to\overline{\mathbb Q}_\ell\xrightarrow{\iota}
\mathbb C$ is injective.  Entrywise injectivity on
general linear groups shows that the complexified
kernel is exactly $\ker\tau_0$.  Thus the integral
type has open kernel, as also recorded in ITA12.

ADP12 now places $F$ in the supremum of the original
$Y_\sigma$, over automorphic irreducible $\sigma$.
For each index with $Y_\sigma=0$, its contribution
is the zero subspace.  For every other index, UPA16
produces an original part $Y$ with $Y.Y=Y_\sigma$
and all the point properties.  Therefore each
$Y_\sigma$ is contained in the supremum over the
parts in (UPA.23), and so is their supremum.
This proves the required membership for $F$.
All initially fixed weight, basis, level, depth,
and function data were arbitrary.  This is the
literal \texttt{AuTeichEx} endpoint.
\end{proof}
% END UPA17

% BEGIN UPA18
\paragraph{UPA18. The bare TJDense endpoint holds for the same original ideal.}
\label{claim:UPA18}
Let $B$ be the original universal ordinary determinant
ring, $\Lambda$ the original weight object, $R_H$ an
original \texttt{UnitaryHeckeData} record, and $K$ an
original permitted level.  Assume the source equality
$V.S'=s.S$.  For every original ideal $J\subset B$,
\[
                  R_H.\operatorname{IsJ}(B,K,J)
                         \quad\Longrightarrow\quad
                  B.\operatorname{Dense}(J).
 \tag{UPA.24}
\]
The original \texttt{HypothesesH'} and all original
coefficient and rank hypotheses are retained; there
is no additional automorphic-Hecke-set premise.
\begin{proof}
The record $R_H$ supplies its literal type
$\tau_0$ and \texttt{TypeForAlt} field.  UPA17
therefore supplies the precise automorphic-Hecke-set
input retained by AFS22.  Apply AFS22 to the original
$B,\Lambda,R_H,K,J$.

For clarity, the conclusion concerns the same ideal
throughout.  In the zero-module branch of
\texttt{IsJ}, one has $J=B$, and the required membership
conclusion is immediate.  In the other branch,
\texttt{IsJ} identifies $J$ with the kernel of the
specified map $f:B\to T$ to the original Hecke algebra.
For an arbitrary increasing weight, AFS1--AFS8 give
its dense good-weight set and the original arithmetic
point dictionary.  AFS20 separates elements of $T$;
AFS21 applies that separation to $f(b)$.
If $b$ vanishes at every original good point killing
$J$, then every separating point of $T$ kills $f(b)$,
so $f(b)=0$ and $b\in\ker f=J$.
This is (UPA.24), with no change of the ideal or
the original ambient ring.
\end{proof}
% END UPA18

\subsection{Scope after installation}

UPA16 supplies the actual Galois-decorated parts and
UPA17 installs them in the original ordinary Hecke setting.
UPA18 completes the previously conditional same-ideal density
endpoint.  The ordinary point output and its flags are developed
in the OFL and HEP chapters, and the full same-ideal
\texttt{RingProps} record in the JRP chapter.  The later all-level
automorphy theorem and both entire rank-four support applications
retain their exact remaining inputs in the current audit.
The conditional irreducible-to-cuspidal assertions above
do not assert that all automorphic representations are
cuspidal or that every Speh multiplicity is one.
